\section{Local walks and entropy}
\label{sec:entropy}

We construct probability distributions on each neighborhood so that rows
indexed by the other two vertices of a triangle are close on average,
with the average weighted by triangle mass.  We obtain the rows by running
a weighted random walk.  The cross-mass estimate bounds the average entropy
by showing that most starting rows, in this weighted average, put almost
all their probability on sets of small weight inside their neighborhoods.
The resulting estimate is independent of the number of vertices.
The use of entropy increments to control averaged distances between
walk distributions has a precedent in Benjamini, Duminil-Copin, Kozma
and Yadin~\cite[Section~2]{BDKY2015}.
We prove the finite weighted version needed here, using the cross-mass
hypothesis to bound its entropy.

Let $G$ be a finite simple graph with positive weights satisfying
\eqref{eq:cross-mass} with constant $C\ge1$.
Fix $x\ge32$ and suppose further that $G$ has an edge and
\begin{equation}
 s_u:=w(N(u))\le e^x\quad\hbox{for every $u$},\qquad
 c_{uv}:=w(N(u)\cap N(v))\ge x^{-1}\quad(uv\in E(G)).
 \label{eq:walk-hypotheses}
\end{equation}
Put
\[
 h_u=\sum_{v\in N(u)}w_vc_{uv}
\]
for every vertex $u$, with the empty sum equal to zero.  We use kernels
only at non-isolated vertices.  At such a vertex $u$, set
\[
 \pi_u(v)=\frac{w_vc_{uv}}{h_u}\quad(v\in N(u)).
\]
These quantities are positive, and
\begin{equation}
 s_u\ge x^{-1},\qquad h_u\ge s_u/x.
 \label{eq:walk-lower-bounds}
\end{equation}
On the finite set $N(u)$ define stochastic matrices
\begin{equation}
 P_u(v,z)=\frac{w_z\mathbf1_{\{vz\in E(G)\}}}{c_{uv}},\qquad
 Q_u=\frac{I+P_u}{2},\qquad K_u^j=P_uQ_u^j\quad(j\ge0).
 \label{eq:walk-kernels}
\end{equation}
Thus $P_u$ is the weighted random walk on the induced graph $G[N(u)]$,
and $Q_u$ adds probability $1/2$ of staying at the current vertex.
The row $K_u^j(v,\cdot)$ is the law after one $P_u$-step from $v$ followed
by $j$ such lazy steps.  The initial step gives
$P_u(v,y)\le xw_y$; we will show that subsequent steps preserve this
density bound.  The holding probability will let us compare rows started
at adjacent vertices through an entropy increment.
The superscript in $K_u^j$ indexes this family, rather than denoting a
power of a fixed kernel.  We shall choose one index $j$ for all neighborhoods.

Here are the probability laws used in the argument.  Put
\[
 Z=\sum_u w_uh_u=6T_G(w)>0,
 \qquad \mu(u,v)=\frac{w_uw_vc_{uv}}{Z}\quad(uv\in E(G)),
\]
where both orientations of each edge occur in the sum defining $Z$.
The conditional distribution of $v$ given $u$ under $\mu$ is $\pi_u$.
If, after sampling $(u,v)$ from $\mu$, we sample $z$ from $P_u(v,\cdot)$,
then
\begin{equation}
 \mathbb P(u,v,z)=\frac{w_uw_vw_z}{Z}
 \quad\hbox{for each ordered triangle $(u,v,z)$}.
 \label{eq:ordered-triangle-law}
\end{equation}
We call this the ordered triangle law.

For a probability vector $a$ supported on vertices with positive weights,
write
\[
 \mathcal H_w(a)=\sum_{y:a(y)>0}a(y)\log\frac{w_y}{a(y)},\qquad
 H_j=\mathbb E_{(u,v)\sim\mu}\mathcal H_w(K_u^j(v,\cdot)).
\]
Thus $H_j$ is an average of row entropies relative to the vertex weights;
the weights themselves need not sum to one.  For probability vectors $a,b$
on the same finite set, write
\[
 \|a-b\|_{\mathrm{TV}}=\frac12\sum_y|a(y)-b(y)|.
\]

The neighborhood bound $s_u\le e^x$ alone allows row entropy as large
as $x$.  The next lemma improves the averaged upper bound to
$O_C(\log x)$ over $\lceil x\rceil$ steps.  Since the entropy increments
are nonnegative, some step has small averaged increment; this will give
the required comparison under the ordered triangle law.

\begin{lemma}[Smoothing the local walks]
\label{lem:entropy-smoothing}
Under \eqref{eq:cross-mass} and \eqref{eq:walk-hypotheses}, put
\[
 E_C=42C^2+2C,\qquad A_C=10+E_C,\qquad B_C=2\sqrt{A_C+1},
 \qquad R=\lceil x\rceil.
\]
Each $K_u^j$ is reversible with stationary measure $\pi_u$, and
\begin{equation}
 K_u^j(v,y)\le xw_y\qquad(j\ge0).
 \label{eq:walk-density}
\end{equation}
Moreover,
\begin{equation}
 -\log x\le H_0\le H_1\le\cdots\le H_R\le A_C\log x.
 \label{eq:entropy-upper}
\end{equation}
There exists $j\in\{0,\ldots,R-1\}$ such that, under the ordered triangle
law,
\begin{equation}
 \mathbb E\bigl\|K_u^j(v,\cdot)-K_u^j(z,\cdot)\bigr\|_{\mathrm{TV}}
 \le B_C\sqrt{\frac{\log x}{x}}.
 \label{eq:triangle-row-distance}
\end{equation}
\end{lemma}

\begin{proof}
We first check the kernels and then bound the entropy they can acquire.
Detailed balance follows from
\[
 \pi_u(v)P_u(v,z)
 =\frac{w_vw_z\mathbf1_{\{vz\in E(G)\}}}{h_u}
 =\pi_u(z)P_u(z,v).
\]
Thus $P_u$ is reversible and stationary for $\pi_u$.  Its powers, and hence
$Q_u$ and $K_u^j$, have the same properties.  No irreducibility assumption
is needed.  By $c_{uv}\ge x^{-1}$, every row of $P_u$ satisfies
$P_u(v,y)\le xw_y$.  Since $P_u$ and $Q_u$ commute,
\[
 K_u^j(v,y)=\sum_zQ_u^j(v,z)P_u(z,y)\le xw_y,
\]
which proves \eqref{eq:walk-density}.  In particular, $H_j\ge-\log x$.

\paragraph{Sets that contain most of the walk.}
Write $L=e^x$.  For each fixed non-isolated starting vertex $v$, we build
sets in the whole graph that contain almost all of the walk for every
center $u\in N(v)$ simultaneously.  Using the same set for all these
centers will allow a cross-mass estimate to control its intersections
with their neighborhoods.  Define
\[
 S_0(v)=N(v),\qquad
 S_{i+1}(v)=S_i(v)\cup
 \{y:w(N(y)\cap S_i(v))\ge x^{-6}/L\}.
\]
We suppress the argument $v$ while it is fixed.  Double counting gives
\[
 \sum_yw_yw(N(y)\cap S_i)
 =\sum_{z\in S_i}w_zs_z\le Lw(S_i).
\]
Consequently
\[
 w(S_{i+1})\le(1+L^2x^6)w(S_i),\qquad
 w(S_R)\le e^x(1+e^{2x}x^6)^R\le e^{20x^2}.
\]
For the last inequality, $x\ge32$ implies $6\log x\le x$,
$\log(1+e^{2x}x^6)\le4x$, and $R\le2x$.

For every edge $uv$ and $0\le i\le R$, we claim that
\begin{equation}
 K_u^i(v,N(u)\setminus S_i(v))\le ix^{-4}.
 \label{eq:walk-leakage}
\end{equation}
The claim holds at $i=0$, since $P_u(v,\cdot)$ is supported on $N(v)$.
In the step from $i$ to $i+1$, mass already outside $S_i$ contributes at
most $ix^{-4}$.  If $z\in S_i$ and $y\notin S_{i+1}$, then $y\ne z$,
so $Q_u(z,y)\le xw_y\mathbf1_{\{zy\in E(G)\}}$.  By
\eqref{eq:walk-density}, the additional mass reaching such a $y\in N(u)$
is at most
\[
 \sum_{z\in N(u)\cap S_i\cap N(y)}xw_z\,xw_y
 \le x^{-4}w_y/L.
\]
Summing over $y\in N(u)$ contributes at most $x^{-4}$, proving the claim.

\paragraph{Most neighborhoods meet these sets in small weight.}
The global set $S_R(v)$ has weight at most $e^{20x^2}$.  We now show
that its intersection with $N(u)$ has weight at most $x^8$ for most
pairs $(u,v)$ under $\mu$.
For fixed $v$ let
\[
 B_v=\{u\in N(v):w(N(u)\cap S_R(v))>x^8\}.
\]
The first use of \eqref{eq:cross-mass}, with parameter $20x^2$,
gives
\[
 x^8w(B_v)\le e_w(N(v),S_R(v))
 \le Cs_v(20x^2+\log^+(1/s_v))
 \le21Cx^2s_v,
\]
where $s_v\ge x^{-1}$ was used in the final inequality.  Thus
$w(B_v)\le21Cx^{-6}s_v$.

This weight estimate must be converted to the conditional law of $u$
given $v$, which is not normalized vertex weight.  The symmetry of $\mu$
gives the exact formula
\[
 \mathbb P_\mu(u\in B_v\mid v)
 =\frac{\sum_{u\in B_v}w_uc_{uv}}{h_v}
 =\frac{e_w(B_v,N(v))}{h_v}.
\]
The second use of \eqref{eq:cross-mass}, now with parameter $x$,
applies because $B_v\subseteq N(v)$ and $s_v\le e^x$.  The function $g_x$
is nondecreasing on $[0,\infty)$, and therefore, for $b\ge0$,
\[
 g_x(b)\le2x\max\{b,e^{-x}\}.
\]
Using $h_v\ge s_v/x$ and $s_v\ge x^{-1}$, we obtain
\begin{align}
 \mathbb P_\mu(u\in B_v\mid v)
 &\le\frac{2Cx^2}{s_v}\bigl(w(B_v)+e^{-x}\bigr)\notag\\
 &\le42C^2x^{-4}+2Cx^3e^{-x}
 \le E_Cx^{-4}.
 \label{eq:exceptional-pair-probability}
\end{align}
Here $e^{-x}\le x^{-7}$ for $x\ge32$.

We can now bound $H_R$.  A probability vector $a$ supported on a set $D$
of total weight $m>0$ satisfies
\begin{equation}
 \mathcal H_w(a)\le\log m.
 \label{eq:weighted-entropy-support}
\end{equation}
Indeed, normalize the weights on $D$ and use nonnegativity of relative
entropy.  If $u\notin B_v$, the set $D=N(u)\cap S_R(v)$ has weight at most
$x^8$, and \eqref{eq:walk-leakage} shows that
$K_u^R(v,D)=1-\delta$ with $0\le\delta\le Rx^{-4}<1$.  Condition the
row on $D$ and its complement.  Formula \eqref{eq:weighted-entropy-support}
and the binary entropy bound then give
\[
 \mathcal H_w(K_u^R(v,\cdot))
 \le\log2+8\log x+\delta x.
\]
An empty complement, or zero probability on the complement, simply
contributes zero.  On pairs with $u\in B_v$, use the unconditional bound
$\mathcal H_w(K_u^R(v,\cdot))\le\log s_u\le x$.  Averaging and applying
\eqref{eq:exceptional-pair-probability} yields
\[
 H_R\le\log2+8\log x+Rx^{-3}+E_Cx^{-3}
 \le(10+E_C)\log x=A_C\log x.
\]

\paragraph{A small entropy increment gives close rows.}
For probability vectors $a,b$, let
\[
 D(a\|b)=\sum_{y:a(y)>0}a(y)\log\frac{a(y)}{b(y)}
\]
denote relative entropy, interpreted as $+\infty$ if necessary.  Sample
$(u,v)$ from $\mu$ and then $z$ from $Q_u(v,\cdot)$.  Since
$Q_uK_u^j=K_u^{j+1}$, the mixture of $K_u^j(z,\cdot)$ over this conditional
choice of $z$ is $K_u^{j+1}(v,\cdot)$.  Consequently, for fixed $u,v$,
\[
 \mathbb E_{z\sim Q_u(v,\cdot)}
 D\bigl(K_u^j(z,\cdot)\|K_u^{j+1}(v,\cdot)\bigr)
 =\mathcal H_w(K_u^{j+1}(v,\cdot))
  -\mathbb E_{z\sim Q_u(v,\cdot)}\mathcal H_w(K_u^j(z,\cdot)).
\]
All these divergences are finite on events of positive probability,
because the second argument is the mixture containing the first.
Conditional on $u$, the law of $v$ is $\pi_u$, which is stationary for
$Q_u$.  After averaging over $v$ and then $u$, the preceding identity
becomes
\begin{equation}
 H_{j+1}-H_j
 =\mathbb E_{\mu Q}
 D\bigl(K_u^j(z,\cdot)\|K_u^{j+1}(v,\cdot)\bigr)\ge0.
 \label{eq:entropy-increment}
\end{equation}
The stationarity step is an averaged identity; it asserts no equality
between individual rows.  Telescoping the increments gives an index
$0\le j<R$ for which
\[
 H_{j+1}-H_j\le\frac{(A_C+1)\log x}{R}
 \le\frac{(A_C+1)\log x}{x}.
\]

For completeness, $\|a-b\|_{\mathrm{TV}}\le\sqrt{D(a\|b)}$ follows
from
\[
 D(a\|b)\ge-2\log\sum_y\sqrt{a(y)b(y)}
 \ge\sum_y(\sqrt{a(y)}-\sqrt{b(y)})^2
\]
and Cauchy--Schwarz.  Therefore \eqref{eq:entropy-increment} and Jensen's
inequality imply
\[
 \mathbb E_{\mu Q}
 \bigl\|K_u^j(z,\cdot)-K_u^{j+1}(v,\cdot)\bigr\|_{\mathrm{TV}}
 \le\sqrt{H_{j+1}-H_j}.
\]
To pass to the ordered triangle law, put
$a_v=K_u^j(v,\cdot)$ and $b_v=K_u^{j+1}(v,\cdot)$ within each fixed
neighborhood.  The triangle inequality and $Q_u=(I+P_u)/2$ give exactly
\begin{align*}
 \mathbb E_{\mu P}\|a_v-a_z\|_{\mathrm{TV}}
 &\le\mathbb E_\mu\|a_v-b_v\|_{\mathrm{TV}}
     +\mathbb E_{\mu P}\|a_z-b_v\|_{\mathrm{TV}}\\
 &=2\mathbb E_{\mu Q}\|a_z-b_v\|_{\mathrm{TV}}.
\end{align*}
Combining the last two bounds proves \eqref{eq:triangle-row-distance}
with $B_C=2\sqrt{A_C+1}$.
\end{proof}
